% !TeX program = lualatex
% Compile twice: lualatex 113.tex
% All references and prose are embedded; no BibTeX database or external inputs.
\documentclass[11pt,a4paper]{article}
\usepackage[margin=24mm,headheight=14pt]{geometry}
\usepackage{fontspec}
\setmainfont{Latin Modern Roman}
\setsansfont{Latin Modern Sans}
\setmonofont{Latin Modern Mono}
\usepackage{amsmath,amssymb,mathtools}
\usepackage{unicode-math}
\setmathfont{Latin Modern Math}
\usepackage{microtype}
\usepackage{enumitem,xurl,needspace}
\usepackage[dvipsnames]{xcolor}
\usepackage{hyperref}
\hypersetup{unicode=true,colorlinks=true,linkcolor=MidnightBlue,urlcolor=MidnightBlue,
 pdftitle={Research attempt: Problem 113},pdfauthor={Research notebook},bookmarksnumbered=true}
\usepackage{bookmark}
\usepackage{tocloft}
\setlength{\cftsecnumwidth}{3em}
\cftsetpnumwidth{2.5em}
\cftsetrmarg{4em}
\usepackage{titlesec}
\titleformat{\section}{\Large\bfseries\raggedright}{\thesection}{0.7em}{}
\usepackage{fancyhdr}
\pagestyle{fancy}\fancyhf{}
\fancyhead[L]{\small\sffamily Open Problems Math | Research notebook}
\fancyhead[R]{\small Problem 113 | 27 September 2026}
\fancyfoot[C]{\thepage}
\setlength{\parindent}{0pt}
\setlength{\parskip}{3pt plus 1pt}
\setlength{\emergencystretch}{3em}
\setcounter{tocdepth}{1}
\setcounter{secnumdepth}{1}
\setlist[itemize]{leftmargin=3em,itemsep=5pt,topsep=4pt,parsep=0pt}
\allowdisplaybreaks
\newcommand{\N}{\mathbb{N}}
\newcommand{\Z}{\mathbb{Z}}
\newcommand{\Q}{\mathbb{Q}}
\newcommand{\R}{\mathbb{R}}
\newcommand{\C}{\mathbb{C}}
\newcommand{\F}{\mathbb{F}}
\newcommand{\A}{\mathbb{A}}
\newcommand{\PP}{\mathbb P}
\newcommand{\E}{\mathbb{E}}
\newcommand{\eps}{\varepsilon}
\newcommand{\dd}{\,\mathrm d}
\newcommand{\sref}[2]{\hyperlink{source:#1:#2}{[S#2]}}
\newcommand{\eref}[2]{\hyperlink{extra:#1:#2}{[E#2]}}
% Turn the next switch off for a shorter reading copy. The quotations preserve
% every exactTarget field, including qualifications omitted from a short summary.
\newif\ifshowcatalogue
\showcataloguetrue
\newcommand{\cataloguescope}[1]{\ifshowcatalogue
 \paragraph{Exact scope in the supplied catalogue.}
 \begin{quote}\small #1\end{quote}\fi}
\usepackage{etoolbox}
\AtBeginEnvironment{itemize}{\small}
\numberwithin{equation}{section}
% Standalone export. Original section 113 preserved verbatim outside marked additions.
\begin{document}
\setcounter{section}{112}
% ================================================================
% problemId: problem.erdos-szemeredi-sum-product-conjecture-for-integers
% source releaseRank: 113; source releaseStatus: open
% exactTarget SHA-256: 2c3ae0da8f5c73969fc2b1484b57a8e4a403ea7b94700d03ef35ac516a390bdf
\clearpage
\section{\texorpdfstring{Erdos-Szemeredi Sum-Product Conjecture for Integers}{Erdos-Szemeredi Sum-Product Conjecture for Integers}}\label{113}
\subsection{Definitions and mathematical statement}
For a finite set $A\subset\Z$, write $A+A=\{a+b:a,b\in A\}$ and $A\cdot A=\{ab:a,b\in A\}$. The conjecture is
\[
 \forall\varepsilon>0\ \exists c_\varepsilon>0\ \forall A\subset\Z\text{ finite}:\qquad
 \max(|A+A|,|A\cdot A|)\ge c_\varepsilon|A|^{2-\varepsilon}.
\]
The constant is uniform in the elements and cardinality of $A$. The assertion says that addition and multiplication cannot both produce exceptionally small image sets. It does not require both sets to have nearly quadratic size.
\par\smallskip\textit{Formulation and concept references:} \sref{113}{1}.
\cataloguescope{Prove that for every epsilon greater than zero there is a constant c\_epsilon greater than zero such that every finite set A of integers satisfies max(|A+A|, |A*A|) at least c\_epsilon times |A|\textasciicircum{}(2-epsilon).}
\subsection{Short English statement}
Must a finite integer set expand almost quadratically under either addition or multiplication?
% BEGIN RESEARCH ADDITION 113
\subsection{Research attempt}

\paragraph{Current frontier and source audit.}
The original source's introduction states the integer conjecture, while
most of that paper investigates sparse-graph variants. We retain the full
sum and product sets here. The real-number version is not interchangeable
with this target. \sref{113}{1}

Cushman's December 2025 manuscript, Theorem~1.3, gives the general-real
exponent $4/3+10/4407-\epsilon$, hence also a lower bound for integer sets.
Bloom--Sawin--Schildkraut--Zhelezov's May 2026 manuscript, Theorem~1.1,
constructs counterexamples over real number fields whose degrees grow
with the set size; its introduction explicitly leaves the integer case
separate. These are identified as research manuscripts, not independently
verified by this notebook. No part of either proof is needed for the
lemmas below. \eref{113}{1}\eref{113}{2}

\paragraph{Attack route.}
A purely real-geometric route must now distinguish itself from the real
counterexample construction. We instead combine exact integer divisibility
with an analysis of how a large translation affects product collisions.
A first height-dependent divisor bound becomes a bound uniform in both
the translation and the step of a containing arithmetic progression.
The remaining issue is to replace that one-dimensional progression
hypothesis by consequences of small additive doubling without destroying
multiplicative relations.

\paragraph{Attempt: the sharp translated-interval threshold.}
For a finite positive set $A$, call its products injective on unordered
pairs if $ab=cd$ with $a,b,c,d\in A$ implies $\{a,b\}=\{c,d\}$ as
multisets; repetitions such as $a=b$ are allowed.

\textbf{Lemma 1.} For integers $D\geq1$ and $M>0$, every subset of
$[M,M+D]\cap\Z$ has injective unordered products whenever
\begin{equation}\label{eq:113:threshold}
 M>\left\lfloor\frac{(D-1)^2}{4}\right\rfloor.
\end{equation}
For every $D\geq3$ the integer threshold is sharp for the full interval.

\emph{Proof.} More generally consider the real translate
$t+B$, where $t>0$ and $B\subseteq\{0,\ldots,D\}$. A collision implies
\[
 t(x+y-z-w)=zw-xy.
\]
If the sums agree, so do the products, and the roots of the corresponding
quadratic show that the unordered pairs agree. Otherwise exchange the
pairs so that $x+y=s$, $z+w=s+\delta$, with integer $\delta\geq1$.
Then $t\delta=xy-zw>0$. The elementary bounds
\[
 xy\leq\frac{s^2}{4},\qquad
 zw\geq\max\{0,D(s+\delta-D)\}
\]
show that no such positive collision is possible when $\delta>D$.
For $1\leq\delta\leq D$, the difference of these bounds is largest at
$s=D-\delta$: below that value it is increasing, and above it its
derivative is $s/2-D\leq0$. Consequently
\[
 t\leq\frac{(D-\delta)^2}{4\delta}
 \leq\frac{(D-1)^2}{4}.
\]
For integer $t=M$ this proves \eqref{eq:113:threshold}.
If $D=2k+1$, the collision at $M=k^2$ is
\[
 (M+k)^2=M(M+2k+1).
\]
If $D=2k$, the collision at $M=k(k-1)$ is
\[
 (M+k-1)(M+k)=M(M+2k).
\]
Both have different unordered pairs, proving sharpness. $\square$

The lemma gives $|AA|=|A|(|A|+1)/2$ in that regime. It also shows why
large absolute height, by itself, is not a plausible source of
counterexamples: a sufficiently remote short interval has maximal
unordered product growth.

\paragraph{Attempt: eliminate the step size by a gcd dichotomy.}
\textbf{Lemma 2.} Fix $C>0$ and $\epsilon>0$. There is a constant
$c_{C,\epsilon}>0$ such that every $n$-element integer set contained in
\[
 a+g\{0,1,\ldots,D\},\qquad g\in\Z\setminus\{0\},\quad
 0\leq D\leq n^C,
\]
satisfies $|AA|\geq c_{C,\epsilon}n^{2-\epsilon}$.
The constant is independent of $a$, $g$, and the elements of $A$.

\emph{Proof.} Discard zero and take either the positive or the negative
part with at least $(n-1)/2$ elements. Negating the latter does not change
its product cardinality. This produces a positive subset $A_0$, with
$m\geq(n-1)/2$, which can be written
\[
 A_0\subseteq M+g\{0,\ldots,D\},\qquad M>0,\quad g>0,
\]
after shrinking and reindexing the containing progression if necessary.
The finitely small cardinalities can be absorbed into the constant.
Divide all elements by $d=\gcd(M,g)$, putting $M'=M/d$, $g'=g/d$;
this preserves product cardinality and gives $\gcd(M',g')=1$.

A product collision among $M'+g'x$ and its companions gives
\begin{equation}\label{eq:113:divisibility}
 M'(x+y-z-w)=g'(zw-xy).
\end{equation}
If $g'>2D$, divisibility forces $x+y-z-w=0$, because its absolute value
is at most $2D$; then the unordered pairs are equal. If
$M'/g'>(D-1)^2/4$, the real-translation version of Lemma~1 also gives
injectivity. In either case $|A_0A_0|=m(m+1)/2$.

In the only remaining case,
\[
 g'\leq2D,\qquad M'\leq g'(D-1)^2/4,
\]
all normalized elements are positive integers at most $3D^3$ for $D\geq1$.
Their products are at most $9D^6$. For completeness, for every $\delta>0$
the divisor function satisfies
\begin{equation}\label{eq:113:divisor}
 d(v)\leq C_\delta v^\delta\qquad(v\geq1).
\end{equation}
Indeed, in $v=\prod p^{e_p}$, for sufficiently large primes
$e+1\leq2^e\leq p^{\delta e}$; for each of the finitely many smaller
primes the supremum of $(e+1)p^{-\delta e}$ is finite. Multiplying these
bounds proves \eqref{eq:113:divisor}.

For each positive product $v$, at most $d(v)$ ordered pairs of positive
integers have that product. Thus
\[
 m^2\leq |A_0A_0|\,C_\delta(9D^6)^\delta
 \leq |A_0A_0|\,C_\delta9^\delta n^{6C\delta}.
\]
Taking $\delta=\epsilon/(6C)$ and using $|AA|\geq|A_0A_0|$ proves the
lemma. The case $D=0$ has at most one element and is harmless. $\square$

This proves a substantial restricted form without using $|A+A|$ at all.
It includes sets of polynomial diameter, but also progressions with
arbitrarily large step and translation. The argument is supplied as a
self-contained deduction, without asserting that the result is new in
the literature.

\paragraph{Attempt to remove the progression hypothesis.}
Small additive doubling often suggests modeling a set by a generalized
arithmetic progression. But neither an arbitrary additive model nor a
translation is a multiplicative homomorphism. The simple divisibility
identity \eqref{eq:113:divisibility} involves only one step; with several
steps, independent linear and quadratic expressions in those steps enter,
and no one-dimensional coprimality argument eliminates all of them.
The precise missing statement for this route is a uniform product-growth
or sum-growth dichotomy for arbitrary integer sets of large normalized
progression span, not a height estimate with an element-dependent constant.

There is also no automatic descent of the 2026 real construction to
integers. For example, if $\alpha$ is an irrational algebraic integer with
irreducible polynomial $f$, there is no unital ring homomorphism
$\Z[\alpha]\to\Z$: the image of $\alpha$ would be an integer root of $f$.
A finite encoding preserving the additions and multiplications expressing
$f(\alpha)=0$ and fixing integers has the same obstruction. Substituting
a large integer for $\alpha$ keeps some additive information but breaks
that exact multiplicative relation. This does not rule out every possible
finite combinatorial transfer; it rules out the naive ring-specialization
step. The growing number-field degree in the manuscript is therefore an
actual target distinction, not just a cosmetic normalization.
\eref{113}{2}

\paragraph{Stress test and exact computations.}
The companion script checks the sharp boundary collisions and the absence
of collisions above the threshold for $3\leq D\leq30$. It also tests the
gcd dichotomy for many progressions, including steps larger than $2D$.
These checks test the formulas; the proofs, not the finite search, establish
the lemmas for every permitted parameter.

The set $A=\{1,2,4,\ldots,2^{n-1}\}$ has $|AA|=2n-1$, so no universal
near-quadratic product-only lemma is possible. Its unordered sums, however,
are all distinct by binary expansion, including the doubled terms, and
$|A+A|=n(n+1)/2$. It is a test against the proposed intermediate assertion,
not against the actual max-of-two-sets conjecture. It also has exponential
span and lies outside the hypotheses just proved.

\paragraph{Outcome.}
\textbf{Outcome: Exploratory approach with unresolved gap.}

Two exact lemmas were obtained: a sharp product-injectivity threshold for
translated integer intervals, and near-quadratic product growth for sets
in polynomial-length arithmetic progressions, uniformly in the progression's
height and step. The remaining arbitrary-set reduction is unproved.
No integer counterexample, full sum-product proof, or improvement to the
unrestricted exponent is claimed.
% END RESEARCH ADDITION 113
\subsection{Sources}
\begin{itemize}\raggedright
\item[\textnormal{[S1]}]\hypertarget{source:113:1}{} Noga Alon, Omer Angel, Itai Benjamini, and Eyal Lubetzky, Sums and products along sparse graphs, Israel Journal of Mathematics 188 (2012), 353-384.
\url{https://www.math.tau.ac.il/~nogaa/PDFS/sumprod2.pdf}.
{\small\textit{Catalogue formulation source.}}
% BEGIN NEW REFERENCES 113
\item[\textnormal{[E1]}]\hypertarget{extra:113:1}{} Adam Cushman,
\emph{A Note on the Sum-Product Problem and the Convex Sumset Problem},
arXiv:2512.13849v1 (15 December 2025), Theorem 1.3.
\url{https://arxiv.org/html/2512.13849v1}.
\item[\textnormal{[E2]}]\hypertarget{extra:113:2}{} Thomas F. Bloom, Will Sawin,
Carl Schildkraut and Dmitrii Zhelezov,
\emph{The sum-product conjecture is false for real numbers},
arXiv:2605.28781v1 (27 May 2026), Theorem 1.1 and Sections 1--2.
The real/algebraic-number target is not the integer target of this entry.
\url{https://arxiv.org/html/2605.28781v1}.
% END NEW REFERENCES 113
\end{itemize}

\end{document}
